\textbf{Phần 1. Trắc nghiệm lựa chọn}
\setcounterpageref{ex}{0}

\begin{ex} %[2D1V02]
Số liền sau của số lớn nhất có hai chữ số là
\choice
{$98$}
{\True $100$}
{$99$}
{$101$}
\loigiai{
Số lớn nhất có hai chữ số là $99$. Số liền sau của 99 là $99 + 1 = 100$.
\textbf{Đáp án của câu hỏi là B}
}
\end{ex}

\begin{ex} %[2D1V02]
Trên tia số, điểm biểu diễn số $A$ cách gốc $0$ là $45$ vạch chia. Điểm $B$ đứng liền sau điểm $A$. Điểm $B$ biểu diễn số
\choice
{$44$}
{$45$}
{$50$}
{\True $46$}
\loigiai{
Điểm $A$ biểu diễn số 45. Điểm liền sau là $45 + 1 = 46$.
\textbf{Đáp án của câu hỏi là D}
}
\end{ex}

\begin{ex} %[2D1V02]
Tìm số tự nhiên $x$, biết số liền trước của $x$ là $58$.
\choice
{\True $59$}
{$57$}
{$60$}
{$58$}
\loigiai{
Vì 58 là số liền trước của $x$ nên $x = 58 + 1 = 59$.
\textbf{Đáp án của câu hỏi là A}
}
\end{ex}

\begin{ex} %[2D1V02]
Biết tổng của số liền trước và số liền sau của số $a$ bằng $40$. Số $a$ là
\choice
{$19$}
{$21$}
{\True $20$}
{$18$}
\loigiai{
Số liền trước là $a - 1$, số liền sau là $a + 1$. Tổng hai số là $(a - 1) + (a + 1) = a + a = 40$. Suy ra $a = 20$.
\textbf{Đáp án của câu hỏi là C}
}
\end{ex}

\begin{ex} %[2D1V02]
Một chú cào cào nhảy trên tia số, mỗi bước nhảy được $5$ đơn vị. Xuất phát từ vạch số $10$, sau $4$ bước nhảy chú cào cào đến vạch số nào?
\choice
{$25$}
{\True $30$}
{$35$}
{$40$}
\loigiai{
Quãng đường cào cào nhảy được là: $4 \times 5 = 20$ đơn vị.\nChú cào cào đến vạch: $10 + 20 = 30$.
\textbf{Đáp án của câu hỏi là B}
}
\end{ex}

\begin{ex} %[2D1V02]
Hai số tự nhiên liên tiếp có tổng bằng $35$. Hai số đó là
\choice
{$16\text{ và } 17$}
{$18\text{ và } 19$}
{$15\text{ và } 20$}
{\True $17\text{ và } 18$}
\loigiai{
Vì hai số liên tiếp hơn kém nhau 1 đơn vị nên số bé là: $(35 - 1) : 2 = 17$. Số lớn là 18.
\textbf{Đáp án của câu hỏi là D}
}
\end{ex}

\begin{ex} %[2D1V02]
Số liền trước của số bé nhất có hai chữ số giống nhau là
\choice
{\True $10$}
{$9$}
{$11$}
{$12$}
\loigiai{
Số bé nhất có hai chữ số giống nhau là $11$. Số liền trước của 11 là $11 - 1 = 10$.
\textbf{Đáp án của câu hỏi là A}
}
\end{ex}

\begin{ex} %[2D1V02]
Trên tia số, khoảng cách giữa số $28$ và số $62$ là bao nhiêu đơn vị?
\choice
{$35\text{ đơn vị}$}
{$44\text{ đơn vị}$}
{\True $34\text{ đơn vị}$}
{$32\text{ đơn vị}$}
\loigiai{
Khoảng cách giữa hai điểm là: $62 - 28 = 34\text{ đơn vị}$.
\textbf{Đáp án của câu hỏi là C}
}
\end{ex}

\begin{ex} %[2D1V02]
Số liền sau của số $x$ là một số tròn chục lớn hơn $40$ nhưng nhỏ hơn $60$. Số $x$ có thể là
\choice
{$59$}
{\True $49$}
{$39$}
{$51$}
\loigiai{
Số tròn chục lớn hơn 40 và nhỏ hơn 60 là 50. Số liền trước của 50 là $50 - 1 = 49$. Do đó $x = 49$.
\textbf{Đáp án của câu hỏi là B}
}
\end{ex}

\begin{ex} %[2D1V02]
Cho ba số tự nhiên liên tiếp tăng dần, biết số ở giữa là $79$. Tích của chữ số hàng đơn vị của số thứ nhất và số thứ ba là
\choice
{$8$}
{$16$}
{$2$}
{\True $0$}
\loigiai{
Ba số liên tiếp là $78, 79, 80$. Số thứ nhất có chữ số hàng đơn vị là 8, số thứ ba có chữ số hàng đơn vị là 0. Tích là $8 \times 0 = 0$.
\textbf{Đáp án của câu hỏi là D}
}
\end{ex}

\begin{ex} %[2D1V02]
Trên tia số có các điểm $A, B, C$. Điểm $A$ là $15$, điểm $B$ là $35$, điểm $C$ nằm chính giữa đoạn $AB$. Điểm $C$ biểu diễn số
\choice
{\True $25$}
{$20$}
{$30$}
{$26$}
\loigiai{
Đoạn $AB$ dài: $35 - 15 = 20$. Nửa đoạn $AB$ là: $20 : 2 = 10$. Điểm $C$ biểu diễn số: $15 + 10 = 25$.
\textbf{Đáp án của câu hỏi là A}
}
\end{ex}

\begin{ex} %[2D1V02]
Tìm số tự nhiên liền trước của số $y$, biết $y + 12 = 45$.
\choice
{$33$}
{$34$}
{\True $32$}
{$31$}
\loigiai{
Ta có: $y = 45 - 12 = 33$. Số liền trước của 33 là $33 - 1 = 32$.
\textbf{Đáp án của câu hỏi là C}
}
\end{ex}

\textbf{Phần 2. Trắc nghiệm đúng sai}
\setcounterpageref{ex}{0}

\begin{ex} %[2D1V02]
Xét tính đúng sai của các khẳng định về số liền trước và số liền sau:
\choiceTF
{\True Hiệu giữa số liền sau và số liền trước của một số bất kì luôn bằng 2.}
{Số 0 có số liền trước là một số tự nhiên.}
{\True Số liền sau của một số chẵn luôn là một số lẻ.}
{\True Ba số tự nhiên liên tiếp luôn có một số chia hết cho 3.}
\loigiai{
\begin{itemize}
    \item Ý a) Đúng: Gọi số đó là $n$, số liền trước là $n-1$, liền sau là $n+1$. Ta có $(n+1) - (n-1) = 2$.
    \item Ý b) Sai: Trong tập hợp các số tự nhiên, số 0 không có số liền trước.
    \item Ý c) Đúng: Một số chẵn cộng thêm 1 đơn vị luôn cho kết quả là một số lẻ.
    \item Ý d) Đúng: Trong ba số tự nhiên liên tiếp chắc chắn có đúng một số chia hết cho 3.
\end{itemize}
\begin{itemize}
    \item a) \textbf{Đúng}.
    \item b) \textbf{Sai}.
    \item c) \textbf{Đúng}.
    \item d) \textbf{Đúng}.
\end{itemize}
}
\end{ex}

\begin{ex} %[2D1V02]
Cho tia số với các mốc cách đều nhau 10 đơn vị từ 0 đến 100:
\choiceTF
{\True Số 47 nằm giữa hai vạch 40 và 50.}
{Số 47 nằm gần vạch 40 hơn vạch 50.}
{\True Từ vạch 20 đến vạch 80 có tất cả 6 khoảng cách 10 đơn vị.}
{Số liền sau vạch 30 là vạch 40.}
\loigiai{
\begin{itemize}
    \item Ý a) Đúng: Vì $40 < 47 < 50$ nên khẳng định đúng.
    \item Ý b) Sai: $47 - 40 = 7$, trong khi $50 - 47 = 3$. Số 47 gần 50 hơn.
    \item Ý c) Đúng: $(80 - 20) : 10 = 6$ khoảng cách.
    \item Ý d) Sai: Số liền sau của 30 là 31, còn 40 là vạch tròn chục tiếp theo.
\end{itemize}
\begin{itemize}
    \item a) \textbf{Đúng}.
    \item b) \textbf{Sai}.
    \item c) \textbf{Đúng}.
    \item d) \textbf{Sai}.
\end{itemize}
}
\end{ex}

\textbf{Phần 3. Trả lời ngắn}
\setcounterpageref{ex}{0}

\begin{ex} %[2D1V02]
Tìm một số biết rằng số liền trước của nó cộng với số liền sau của nó bằng 84.
\shortans{42}
\loigiai{
\textbf{Đáp số:} $42$\\
Tổng số liền trước và liền sau gấp đôi số đó: $84 : 2 = 42$.
}
\end{ex}

\begin{ex} %[2D1V02]
Số liền trước của số lớn nhất có hai chữ số khác nhau là số nào?
\shortans{97}
\loigiai{
\textbf{Đáp số:} $97$\\
Số lớn nhất có hai chữ số khác nhau là 98. Số liền trước là 97.
}
\end{ex}

\begin{ex} %[2D1V02]
Một chú thỏ nhảy từ vạch 12 trên tia số, mỗi lần nhảy 3 bước. Sau bao nhiêu lần nhảy thì thỏ đến vạch 33?
\shortans{7}
\loigiai{
\textbf{Đáp số:} $7$\\
Khoảng cách thỏ nhảy: $33 - 12 = 21$. Số lần nhảy là: $21 : 3 = 7$ lần.
}
\end{ex}

\begin{ex} %[2D1V02]
Tìm số tự nhiên biết rằng nếu cộng số đó với số liền sau của nó thì được 51.
\shortans{25}
\loigiai{
\textbf{Đáp số:} $25$\\
Hai số tự nhiên liên tiếp có tổng là 51: số bé là $(51 - 1) : 2 = 25$.
}
\end{ex}

\begin{ex} %[2D1V02]
Trên tia số, giữa số 18 và số 42 có tất cả bao nhiêu số tự nhiên?
\shortans{23}
\loigiai{
\textbf{Đáp số:} $23$\\
Số các số nằm ở giữa: $(42 - 18) - 1 = 23$ số.
}
\end{ex}

\begin{ex} %[2D1V02]
Tìm số liền sau của số tròn chục lớn nhất có hai chữ số mà tổng hai chữ số bằng 9.
\shortans{91}
\loigiai{
\textbf{Đáp số:} $91$\\
Số tròn chục có dạng $\overline{a0} \Rightarrow a + 0 = 9 \Rightarrow a = 9$. Số đó là 90, số liền sau là 91.
}
\end{ex}
